% Created 2026-03-14 Sat 06:02
% Intended LaTeX compiler: pdflatex
\documentclass[11pt]{article}
\usepackage[utf8]{inputenc}
\usepackage[T1]{fontenc}
\usepackage{graphicx}
\usepackage{longtable}
\usepackage{wrapfig}
\usepackage{rotating}
\usepackage[normalem]{ulem}
\usepackage{amsmath}
\usepackage{amssymb}
\usepackage{capt-of}
\usepackage{hyperref}
\author{The Prologos Project - from Zachary Larson}
\date{2026-03-14}
\title{Formal Modeling on Propagator Networks\\\medskip
\large The Landscape of Temporal Logic, Fixpoint Theory, and Formal Verification --- and How Propagators Embed It All}
\hypersetup{
 pdfauthor={The Prologos Project - from Zachary Larson},
 pdftitle={Formal Modeling on Propagator Networks},
 pdfkeywords={},
 pdfsubject={},
 pdfcreator={Emacs 30.2 (Org mode 9.7.39)}, 
 pdflang={English}}
\begin{document}

\maketitle
\tableofcontents

\section{Abstract}
\label{sec:orged3eafc}

This document surveys the landscape of formal modeling, temporal logic, and
verification, with a sustained focus on how propagator networks serve as a
universal computational substrate for these formalisms. The core observation,
developed in our earlier research (\texttt{PROPAGATORS\_AS\_MODEL\_CHECKERS.org}), is
that model checking, abstract interpretation, and temporal reasoning are all
\emph{monotone fixpoint computations over complete lattices} --- and propagator
networks are precisely the engine for such computations.

We cast a wide net: from the modal \(\mu\)-calculus and its alternation hierarchy
through signal temporal logic and continuous-domain verification; from
Approximation Fixpoint Theory and bilattices through coinduction and greatest
fixpoints; from runtime verification and shield synthesis through reactive
synthesis and separation logic. In each case, we identify the lattice-theoretic
core and trace the connection to propagator network computation. The document
concludes with a synthesis of how these formalisms map to Prologos's existing
and planned infrastructure, and a research program for bringing formal
verification from expert-only tools to type-system-level accessibility.

This is a stage-0 research document: it maps the territory and identifies
implementation tracks, but does not commit to specific implementation timelines.
\section{1. The Modal \(\mu\)-Calculus}
\label{sec:orgaefd4df}

\subsection{1.1 Foundations}
\label{sec:orga8cd2ad}

The modal \(\mu\)-calculus (Kozen 1983) is the most expressive logic in the
temporal/modal family that retains decidable model checking. It extends
propositional modal logic with explicit least (\(\mu\)) and greatest (\(\nu\)) fixpoint
operators:

\(\phi\) ::= p | \(\neg{}\) \(\phi\) | \(\phi\) \(\land\) \(\psi\) | \(\phi\) \(\lor\) \(\psi\)
\begin{center}
\begin{tabular}{ll}
\(\langle\) a \(\rangle\) \(\phi\) & {[}a] \(\phi\)\\
X & \(\mu\) X. \(\phi\)(X) & \(\nu\) X. \(\phi\)(X)\\
\end{tabular}
\end{center}

where X occurs positively in \(\phi\)(X) (under an even number of negations).

\textbf{Semantics}: Formulas are interpreted over labeled transition systems
(S, \{R\textsubscript{a}\}\textsubscript{a \(\in\) Act}, L). The denotation \llbracket \(\phi\) \rrbracket \subseteq S
is a subset of states. The fixpoint operators have set-theoretic semantics via
the Knaster-Tarski theorem:

\begin{itemize}
\item \llbracket \(\mu\) X. \(\phi\)(X) \rrbracket = \bigcap $\backslash${ T \subseteq S \mid F(T) \subseteq T $\backslash$}
(least pre-fixpoint = least fixpoint for monotone F)
\item \llbracket \(\nu\) X. \(\phi\)(X) \rrbracket = \bigcup $\backslash${ T \subseteq S \mid T \subseteq F(T) $\backslash$}
(greatest post-fixpoint = greatest fixpoint)
\end{itemize}

where F(T) = \llbracket \(\phi\)(X := T) \rrbracket is monotone by the positivity
requirement.

When the state space is finite, fixpoints are computable by Kleene iteration:
\(\mu\) X. \(\phi\)(X) = \bigcup\textsubscript{i \(\ge\) 0} F\textsuperscript{i}(\(\emptyset\)), converging in at most |S| steps.
Dually, \(\nu\) X. \(\phi\)(X) = \bigcap\textsubscript{i \(\ge\) 0} F\textsuperscript{i}(S).

\textbf{Key references}:
\begin{itemize}
\item Kozen, D. "Results on the propositional \(\mu\)-calculus." \emph{TCS} 27(3), 1983.
\item Walukiewicz, I. "Completeness of Kozen's axiomatisation." \emph{Information and Computation} 157(1-2), 2000.
\end{itemize}
\subsection{1.2 Expressiveness: The Janin-Walukiewicz Theorem}
\label{sec:org1fd1d4a}

The \(\mu\)-calculus has exactly the right expressiveness for model checking:

\textbf{Theorem} (Janin-Walukiewicz 1996): The modal \(\mu\)-calculus is expressively
equivalent to the bisimulation-invariant fragment of monadic second-order logic
(MSO) over trees.

This places the \(\mu\)-calculus at a precise point in the logical landscape: it
captures everything that MSO can say about transition systems, up to
bisimulation equivalence. Since bisimulation is the natural equivalence for
concurrent systems, the \(\mu\)-calculus is the "right" logic for process
verification.
\subsection{1.3 The Alternation Hierarchy}
\label{sec:orge4900ce}

The alternation depth of a \(\mu\)-calculus formula measures the nesting depth of
mutually recursive least and greatest fixpoint operators. Bradfield (1996) and
Lenzi (1996) independently proved:

\textbf{Theorem} (Bradfield 1996): The modal \(\mu\)-calculus alternation hierarchy is
\textbf{strictly infinite}. For every n, there exist properties expressible at
alternation depth n+1 but not at depth n.

The hierarchy levels correspond to increasing complexity of temporal properties:

\begin{center}
\begin{tabular}{rlll}
Level & \(\mu\)-Calculus Pattern & Temporal Meaning & Example\\
\hline
0 & p (no fixpoint) & State property & "this state is labeled p"\\
1 & \(\mu\) X. \(\phi\) or \(\nu\) X. \(\phi\) & Safety/reachability & AG p, EF p\\
2 & \(\nu\) X. \(\mu\) Y. \(\phi\)(X,Y) & Response/persistence & AG(p \(\to\) AF q), AGEF p\\
3+ & Deeper alternation & Fairness, complex liveness & Streett/Rabin conditions\\
\end{tabular}
\end{center}

\textbf{For Prologos}: Our current session type checker operates at Level 1 (safety
via contradiction detection, reachability via completeness checking). Level 2
properties (response: "every request gets a reply") require nested fixpoints
--- computationally feasible in propagator networks but needing an interface
for specification. Level 3+ properties appear rarely in session type contexts
but are needed for fairness in multi-party protocols.
\subsection{1.4 Parity Games: The Decision Procedure}
\label{sec:org2348743}

Model checking \(\mu\)-calculus formulas reduces to solving \emph{parity games} --- two-player
infinite games where vertices are colored with natural numbers (priorities) and
the winner is determined by the parity (even/odd) of the highest priority seen
infinitely often.

\textbf{Theorem} (Emerson-Jutla-Sistla 1993): The \(\mu\)-calculus model checking problem
is equivalent to solving parity games.

The complexity of parity games was a major open problem for decades:
\begin{itemize}
\item The problem lies in NP \(\cap\) coNP (Emerson-Jutla-Sistla 1993)
\item It is in UP \(\cap\) coUP (Jurdzinski 1998)
\item \textbf{Quasipolynomial-time algorithms} were discovered by Calude, Jain, Khoussainov,
Li, and Stephan (STOC 2017), resolving a long-standing conjecture. The algorithm
runs in time n\textsuperscript{O(\(\log\) d)} where n is the game size and d is the number of
priorities.
\item Jurdzinski and Lazic (LICS 2017) gave a succinct progress measure algorithm
also achieving quasipolynomial time.
\end{itemize}

\textbf{For propagators}: Parity games correspond to mixed lfp/gfp computations with
alternating fixpoints. The game-theoretic view suggests that propagator networks
computing nested fixpoints can be understood as strategies in parity games,
with the scheduler playing the role of the game solver.
\subsection{1.5 The Automata-Logic-Games Triangle}
\label{sec:org7039cac}

A deep structural equivalence connects three formalisms:

\begin{verbatim}
  Alternating Parity Tree Automata
            /                   \
           /                     \
Modal μ-Calculus  ←————→   Parity Games
\end{verbatim}

\begin{itemize}
\item \(\mu\)-calculus formulas translate to alternating parity tree automata (Wilke 2001)
\item Emptiness of alternating parity tree automata reduces to parity games
\item Parity games can be encoded as \(\mu\)-calculus model checking problems
\end{itemize}

This triangle means results in any one formalism transfer to the others. The
quasipolynomial breakthrough for parity games immediately gives
quasipolynomial \(\mu\)-calculus model checking and automata emptiness.
\subsection{1.6 Extensions}
\label{sec:orge6b1bb3}

\subsubsection{Coalgebraic \(\mu\)-Calculus}
\label{sec:org498ae5f}

Cirstea, Kupke, and Pattinson generalize the \(\mu\)-calculus from Kripke frames to
coalgebras for arbitrary endofunctors F : Set \(\to\) Set. The key innovation is
\emph{predicate liftings} \(\lambda\) : Q\textsuperscript{n} \(\to\) Q \^{} F that replace the modal operators
\(\langle\) a \(\rangle\) and [a] with functorial generalizations. This captures:
\begin{itemize}
\item Probabilistic modal logic (F = distribution functor)
\item Graded modal logic (F = multiset/bag functor)
\item Coalition logic (F = coalition functor)
\end{itemize}

The coalgebraic perspective is categorically dual to inductive definitions ---
where least fixpoints are initial algebras, greatest fixpoints are terminal
coalgebras. This connects directly to the coinductive reasoning Prologos needs
for recursive session types.
\subsubsection{Quantitative and Vectorial \(\mu\)-Calculus}
\label{sec:org141d43f}

Fischer and Grädel (LICS 2023) develop a \emph{vectorial} \(\mu\)-calculus where fixpoint
variables range over tuples of sets, enabling systems of mutually recursive
fixpoint equations. This corresponds to multi-cell propagator interactions
where cells mutually constrain each other.

Lattice-valued \(\mu\)-calculus replaces the Boolean powerset 2\textsuperscript{S} with an arbitrary
complete lattice L, yielding L-valued satisfaction. This is precisely the
generalization that propagator networks already implement: cells hold lattice
values, not Boolean sets.
\subsubsection{Fixpoint Logic with Chop (FLC)}
\label{sec:orge0947d5}

Müller-Olm (1999) extends the \(\mu\)-calculus with a sequential composition operator
("chop"), enabling specification of properties that depend on the \emph{sequential
structure} of computations. This is relevant for session types, which are
inherently sequential protocols.
\section{2. The Temporal Logic Landscape}
\label{sec:org1be338f}

\subsection{2.1 The Classical Core: LTL, CTL, CTL*}
\label{sec:orgd781ad9}

Three fundamental temporal logics form the core of the landscape:

\textbf{LTL} (Pnueli 1977) reasons about individual computation paths with operators
X (next), U (until), F (eventually), G (always). It captures exactly the
star-free \(\omega\)-regular languages (Kamp's theorem). Model checking is
PSPACE-complete.

\textbf{CTL} (Clarke-Emerson 1981) reasons about branching computation trees. Every
temporal operator is paired with a path quantifier (A = all paths, E = exists
path). Model checking is polynomial: O(|\(\phi\)| \(\cdot\) |M|) via iterative fixpoint
computation on the powerset lattice.

\textbf{CTL*} (Emerson-Halpern 1986) subsumes both by allowing free mixing of path
quantifiers and temporal operators. CTL and LTL are \emph{incomparable} fragments:
CTL can express existential path properties (EG p) that LTL cannot; LTL can
express fairness properties (FG p \(\to\) GF q) that CTL cannot.

All three embed into the modal \(\mu\)-calculus:

\begin{center}
\begin{tabular}{ll}
Temporal Logic & \(\mu\)-Calculus Translation\\
\hline
AG \(\phi\) & \(\nu\) X. \(\phi\) \(\land\) {[}a]X\\
EF \(\phi\) & \(\mu\) X. \(\phi\) \(\lor\) \(\langle\) a \(\rangle\) X\\
AF \(\phi\) & \(\mu\) X. \(\phi\) \(\lor\) {[}a]X\\
EG \(\phi\) & \(\nu\) X. \(\phi\) \(\land\) \(\langle\) a \(\rangle\) X\\
A(\(\phi\) U \(\psi\)) & \(\mu\) X. \(\psi\) \(\lor\) (\(\phi\) \(\land\) {[}a]X \(\land\) \(\langle\) a \(\rangle\) \top)\\
E(\(\phi\) U \(\psi\)) & \(\mu\) X. \(\psi\) \(\lor\) (\(\phi\) \(\land\) \(\langle\) a \(\rangle\) X)\\
\end{tabular}
\end{center}

CTL embeds into the alternation-free fragment (Level 1). CTL* requires
alternation depth 2. The Emerson-Lei algorithm computes nested fixpoints in
O(n\textsuperscript{d}) symbolic steps where d is the alternation depth.
\subsection{2.2 Signal Temporal Logic (STL)}
\label{sec:orgf151770}

STL (Maler-Nickovic 2004; Fainekos-Pappas 2009; Donzé-Maler 2010) brings
temporal logic to continuous-time, real-valued signals --- the natural output
of cyber-physical systems.

\textbf{Syntax}:
\(\phi\) ::= \(\mu\) | \(\neg{}\) \(\phi\) | \(\phi\) \(\land\) \(\psi\) | F\textsubscript{{[}a,b]} \(\phi\) | G\textsubscript{{[}a,b]} \(\phi\) | \(\phi\) U\textsubscript{{[}a,b]} \(\psi\)

where \(\mu\) is a predicate f(s(t)) > 0 over signal values. All temporal operators
carry explicit time bounds [a,b].

\textbf{The Robustness Innovation}: The key contribution of STL is \emph{quantitative
robustness semantics}. Instead of Boolean satisfaction, the robustness degree
\(\rho\)(\(\phi\), s, t) \(\in\) \overline\{\mathbb{R}\} measures \emph{how robustly} a signal
satisfies or violates a specification:

\begin{itemize}
\item \(\rho\)(\(\mu\), s, t) = f(s(t)) (raw predicate value)
\item \(\rho\)(\(\neg{}\) \(\phi\), s, t) = -\(\rho\)(\(\phi\), s, t)
\item \(\rho\)(\(\phi\) \(\land\) \(\psi\), s, t) = \(\min\)(\(\rho\)(\(\phi\), s, t), \(\rho\)(\(\psi\), s, t))
\item \(\rho\)(\(\phi\) \(\lor\) \(\psi\), s, t) = \(\max\)(\(\rho\)(\(\phi\), s, t), \(\rho\)(\(\psi\), s, t))
\item \(\rho\)(G\textsubscript{{[}a,b]} \(\phi\), s, t) = \(\inf\)\textsubscript{t' \(\in\) {[}t+a, t+b]} \(\rho\)(\(\phi\), s, t')
\item \(\rho\)(F\textsubscript{{[}a,b]} \(\phi\), s, t) = \(\supset\)\textsubscript{t' \(\in\) {[}t+a, t+b]} \(\rho\)(\(\phi\), s, t')
\end{itemize}

\textbf{Soundness}: \(\rho\) > 0 implies satisfaction; \(\rho\) < 0 implies violation. The
magnitude |\(\rho\)| quantifies the margin.

\textbf{Lattice-theoretic structure}: STL replaces the Boolean lattice \{false, true\}
with the complete lattice (\overline\{\mathbb{R}\}, \(\le\)) where min/max replace
meet/join, and inf/sup over time intervals replace the fixpoint iteration.
This is exactly the pattern propagator networks compute: cells hold lattice
values, propagators are monotone functions, and run-to-quiescence computes
the fixpoint.

\textbf{Applications}: Online/offline monitoring of CPS, robustness-guided
falsification for automotive/avionics testing, control synthesis via
optimization-based MPC, and differentiable STL for gradient-based learning.

\textbf{Key references}:
\begin{itemize}
\item Maler, O. and Nickovic, D. "Monitoring temporal properties of continuous signals." \emph{FORMATS} 2004.
\item Fainekos, G. and Pappas, G. "Robustness of temporal logic specifications for continuous-time signals." \emph{TCS} 410(42), 2009.
\item Donzé, A. and Maler, O. "Robust satisfaction of temporal logic over real-valued signals." \emph{FORMATS} 2010.
\end{itemize}
\subsection{2.3 Metric Temporal Logic (MTL)}
\label{sec:orge631957}

MTL (Koymans 1990) extends LTL with time-bounded temporal operators interpreted
over \emph{timed words} --- sequences of events with real-valued timestamps.

\(\phi\) ::= p | \(\neg{}\) \(\phi\) | \(\phi\) \(\land\) \(\psi\) | \(\phi\) U\textsubscript{I} \(\psi\)

where I is an interval (e.g., [0,5], (3, \(\infty\))). The decidability landscape
is surprisingly delicate:

\begin{center}
\begin{tabular}{llll}
Domain & Problem & Decidability & Complexity\\
\hline
Finite timed words & Satisfiability & Decidable & Non-primitive recursive\\
Infinite timed words & Satisfiability & Undecidable & ---\\
Finite timed words & Model checking & Decidable & Non-primitive recursive\\
Infinite timed words & Model checking & Undecidable & ---\\
\end{tabular}
\end{center}

The decidability of MTL over finite timed words was established by Ouaknine and
Worrell (LICS 2005), resolving a longstanding open question.

\textbf{MITL} (Alur-Feder-Henzinger 1996) bans punctual constraints (exact-time like
U\textsubscript{=5}), achieving EXPSPACE-completeness. The slogan: "the cost of punctuality"
--- exact timing is the source of undecidability.

\textbf{Connection to timed automata}: MTL formulas can be translated to one-clock
alternating timed automata (OCATA), which bridge to timed automata (Alur-Dill
\begin{enumerate}
\item for model checking.
\end{enumerate}
\subsection{2.4 HyperLTL and Hyperproperties}
\label{sec:org8d9c44e}

Clarkson and Schneider (JCS 2010) introduced \emph{hyperproperties} --- properties
of \emph{sets of traces} rather than individual traces. This generalization is
necessary because security properties like non-interference and opacity relate
multiple execution traces:

\textbf{HyperLTL} (Clarkson-Finkbeiner-Koleini-Rabe-Sanchez, POST 2014) adds trace
quantifiers to LTL:

\(\psi\) ::= \(\forall\) \(\pi\). \(\psi\) | \(\exists\) \(\pi\). \(\psi\) | \(\phi\)

where \(\phi\) is an LTL formula with indexed atomic propositions a\textsubscript{\(\pi\)}.

\textbf{Expressible properties}:
\begin{itemize}
\item \emph{Non-interference}: \(\forall\) \(\pi\). \(\exists\) \(\pi\)'. (lo\textsubscript{\(\pi\)} = lo\textsubscript{\(\pi\)'}) \(\land\) G(out\textsubscript{\(\pi\)} \(\leftrightarrow\) out\textsubscript{\(\pi\)'})
\item \emph{Observational determinism}: \(\forall\) \(\pi\). \(\forall\) \(\pi\)'. G(l\textsubscript{\(\pi\)} \(\leftrightarrow\) l\textsubscript{\(\pi\)'}) \(\to\) G(o\textsubscript{\(\pi\)} \(\leftrightarrow\) o\textsubscript{\(\pi\)'})
\item \emph{Opacity}: \(\forall\) \(\pi\). (secret\textsubscript{\(\pi\)} \(\to\) \(\exists\) \(\pi\)'. \(\neg{}\) secret\textsubscript{\(\pi\)'} \(\land\) G(obs\textsubscript{\(\pi\)} \(\leftrightarrow\) obs\textsubscript{\(\pi\)'}))
\end{itemize}

Model checking is decidable for the alternation-free fragment but undecidable
with quantifier alternation (\(\forall\) \(\exists\)). Monitoring alternation-free HyperLTL
is feasible; monitoring with alternation is fundamentally limited.

\textbf{For Prologos}: HyperLTL's trace quantifiers mirror the ATMS worldview
structure. Universal hyperproperties (\(\forall\) \(\pi\)) correspond to "all worldviews
satisfy"; existential (\(\exists\) \(\pi\)) to "some worldview satisfies." This
connection suggests that Prologos could verify hypersecurity properties of
session-typed protocols using existing ATMS infrastructure.
\subsection{2.5 Probabilistic Temporal Logics}
\label{sec:org574f917}

\textbf{PCTL} (Hansson-Jonsson 1994) replaces CTL's path quantifiers A/E with
probability thresholds P\textsubscript{\(\ge\) p}. Model checking over DTMCs reduces to
solving linear equation systems. Over MDPs, it asks whether there exists a
policy achieving the probability bound, solvable via linear programming.

\textbf{CSL} (Baier-Haverkort-Hermanns-Katoen 2003) extends PCTL for continuous-time
Markov chains with time-bounded until and steady-state operators. Model
checking uses uniformisation for transient probability computation.

\textbf{For propagators}: Probabilistic verification replaces the Boolean lattice with
{[}0,1]. Propagator cells holding probability values, refined by probabilistic
transition propagators, would compute the same fixpoints as PRISM/Storm.
Widening operators for probability lattices (interval abstraction) could
accelerate convergence.
\subsection{2.6 Spatial and Spatio-Temporal Logics}
\label{sec:orga8a1ae0}

Caires and Cardelli (2001-2004) developed a \emph{spatial logic for concurrency}
with a parallel composition operator and its adjoint (guarantee), mirroring
multiplicative linear logic:

\begin{itemize}
\item \(\phi\) | \(\psi\): parallel composition (separating conjunction, like * in
separation logic)
\item \(\phi\) \triangleright \(\psi\): guarantee (linear implication)
\end{itemize}

This connects temporal logic directly to \emph{resource reasoning}: spatial
reasoning about concurrent processes is fundamentally resource-sensitive.
The ambient logic (Cardelli-Gordon 2000) extends this to mobile computation
with nested locations.

Topological semantics for modal logic (McKinsey-Tarski 1944) interprets
\Box \(\phi\) as the interior and \(\diamondsuit\) \(\phi\) as the closure, identifying
S4 modal logic with topological spaces.
\subsection{2.7 Quantitative and Weighted Temporal Logics}
\label{sec:orgdb916fe}

\textbf{Discounted LTL} (Almagor-Boker-Kupferman, TACAS 2014) replaces Boolean
satisfaction with [0,1]-valued satisfaction, discounting temporal distance:
\llbracket F \(\phi\) \rrbracket(w,i) = \(\supset\)\textsubscript{j \(\ge\) i} \(\lambda\)\textsuperscript{j-i} \(\cdot\) \llbracket \(\phi\) \rrbracket(w,j).

\textbf{Energy/mean-payoff games} (Chatterjee-Doyen-Henzinger 2010) combine parity
conditions with quantitative resource constraints. Energy parity games
(ensuring resource never drops below zero while satisfying a \(\omega\)-regular
condition) are decidable in NP \(\cap\) coNP.

\textbf{For propagators}: Quantitative temporal logics correspond to propagator
computations over richer lattices than the Boolean powerset. The energy lattice
(\mathbb{Z}, \(\le\)) with min as meet, the discount lattice ([0,1], \(\le\)) with
sup/inf --- these are all complete lattices amenable to propagator computation.
The key insight is that the \emph{same propagator scheduler} handles all lattices;
only the cell merge function changes.
\section{3. Fixpoint Theory and Propagator Networks}
\label{sec:org13c6079}

\subsection{3.1 The Knaster-Tarski Foundation}
\label{sec:orgc0731d4}

The Knaster-Tarski theorem is the mathematical bedrock of both model checking
and propagator computation:

\textbf{Theorem} (Knaster 1928, Tarski 1955): Let L be a complete lattice and
f : L \(\to\) L a monotone function. Then the set of fixpoints of f forms a complete
lattice, with:
\begin{itemize}
\item lfp(f) = \bigcap $\backslash${ x \(\in\) L \mid f(x) \(\le\) x $\backslash$} (meet of pre-fixpoints)
\item gfp(f) = \bigcup $\backslash${ x \(\in\) L \mid x \(\le\) f(x) $\backslash$} (join of post-fixpoints)
\end{itemize}

\textbf{Kleene's theorem} strengthens this for \(\omega\)-continuous functions (those
preserving directed suprema): lfp(f) = \bigsqcup\textsubscript{i \(\ge\) 0} f\textsuperscript{i}(\bot), computed
by iterating from bottom. Dually, gfp(f) = \bigsqcap\textsubscript{i \(\ge\) 0} f\textsuperscript{i}(\top).

Propagator networks are \emph{precisely} Kleene iteration engines:
\begin{itemize}
\item Each cell starts at \bot (ascending) or \top (descending)
\item Each propagator is a monotone function
\item \texttt{run-to-quiescence} iterates until no cell changes --- the fixpoint
\end{itemize}
\subsection{3.2 Approximation Fixpoint Theory (AFT)}
\label{sec:orgb4e38ee}

Denecker, Marek, and Truszczyński (2000-2004) developed Approximation Fixpoint
Theory to unify the semantics of non-monotone reasoning formalisms. The key
construction:

Given a lattice L, form the \emph{bilattice} L\textsuperscript{2} = L \texttimes{} L with two orderings:
\begin{itemize}
\item \emph{Precision ordering} \(\le\)\textsubscript{p}: (x\textsubscript{1}, x\textsubscript{2}) \(\le\)\textsubscript{p} (y\textsubscript{1}, y\textsubscript{2}) iff x\textsubscript{1} \(\le\) y\textsubscript{1} and y\textsubscript{2} \(\le\) x\textsubscript{2}
(more precise = tighter approximation)
\item \emph{Truth ordering} \(\le\)\textsubscript{t}: (x\textsubscript{1}, x\textsubscript{2}) \(\le\)\textsubscript{t} (y\textsubscript{1}, y\textsubscript{2}) iff x\textsubscript{1} \(\le\) y\textsubscript{1} and x\textsubscript{2} \(\le\) y\textsubscript{2}
(more true = higher in both components)
\end{itemize}

A pair (x, y) \(\in\) L\textsuperscript{2} represents an \emph{approximation}: x is the lower bound
(what is certainly true) and y is the upper bound (what is possibly true). An
exact element has x = y.

\textbf{Stable operators}: An operator A : L\textsuperscript{2} \(\to\) L\textsuperscript{2} is an approximation of an
operator O : L \(\to\) L if for exact (x,x), A(x,x) = (O(x), O(x)). AFT defines:
\begin{itemize}
\item \emph{Kripke-Kleene fixpoint}: lfp\textsubscript{\(\le\)\textsubscript{p}}(A) --- the most precise fixpoint in the
precision ordering
\item \emph{Stable fixpoint}: a fixpoint that is minimal in a certain well-founded
iteration
\item \emph{Well-founded fixpoint}: always exists and is unique
\end{itemize}

\textbf{Unification power}: AFT simultaneously captures:
\begin{itemize}
\item Stable model semantics (answer set programming)
\item Well-founded semantics (logic programming with negation)
\item Default logic fixpoints
\item Autoepistemic logic expansions
\end{itemize}

\textbf{For Prologos}: The bilattice L\textsuperscript{2} maps directly to a two-cell propagator
pattern: one cell for the lower bound (ascending), one for the upper bound
(descending). The precision ordering is the product lattice ordering. AFT's
approximation operators are exactly cross-cell propagators that enforce the
consistency between lower and upper bounds. This gives Prologos a principled
way to handle non-monotone operations (like negation-as-failure in the logic
engine) within a monotone propagator framework --- which is precisely what
the stratification controller already does, but without the theoretical
framework to generalize it.
\subsection{3.3 Belnap's Four-Valued Logic and Bilattices}
\label{sec:org4eca73b}

Belnap (1977) introduced a four-valued logic FOUR = \{\bot, t, f, \top\}
for reasoning with incomplete and contradictory information:

\begin{verbatim}
  ⊤ (both/contradictory)
 / \
t   f
 \ /
  ⊥ (neither/unknown)
\end{verbatim}

This carries two orderings:
\begin{itemize}
\item \emph{Knowledge} (information): \bot < t,f < \top
\item \emph{Truth}: f < \bot, \top < t
\end{itemize}

Ginsberg (1988) generalized to \emph{bilattices} --- structures with two lattice
orderings. Fitting (1991-2002) applied bilattices to logic programming,
interpreting programs over FOUR to handle incomplete and contradictory
information simultaneously.

\textbf{For Prologos}: The session lattice \texttt{sess-bot} / concrete / \texttt{sess-top} is
already a three-valued approximation of FOUR. The ATMS adds a fourth value
(contradictory worldview). The connection to bilattice logic programming
suggests a path toward handling contradictory constraints gracefully rather
than treating contradiction as a hard failure.
\subsection{3.4 Mixed Least/Greatest Fixpoints}
\label{sec:orgb43798b}

The central challenge of the \(\mu\)-calculus is interleaving least and greatest
fixpoints. A formula \(\nu\) X. \(\mu\) Y. \(\phi\)(X, Y) requires computing:
\begin{enumerate}
\item For a given approximation of X, compute the least fixpoint of Y \mapsto \(\phi\)(X, Y)
\item Use the result to refine X toward its greatest fixpoint
\item Iterate until both X and Y stabilize
\end{enumerate}

The \emph{parity condition} is the general pattern for such interleaving: assign
a priority (natural number) to each fixpoint variable, and determine the
winner by the parity of the highest priority seen infinitely often. This is
why parity games are the decision procedure for the \(\mu\)-calculus.

\textbf{Progress measures} (Jurdzinski 2000) provide an alternative characterization:
a progress measure is a function that assigns to each state a value from a
well-founded set, such that the function decreases along certain edges.
Computing the least progress measure solves the parity game.

\textbf{For propagator networks}: Mixed lfp/gfp requires cells with \emph{different
directions} in the same network:
\begin{itemize}
\item Ascending cells (lfp): start at \bot, refine up via join
\item Descending cells (gfp): start at \top, refine down via meet
\end{itemize}

The scheduler must handle both directions, with nested fixpoints computed
inside-out. This is architecturally straightforward --- the current BSP/Jacobi
scheduler already processes all cells in rounds; adding descending cells
requires only changing the merge operation per cell. The deeper question is
termination for mixed networks, addressed by Baldan et al.'s fixpoint games
(§3.7).
\subsection{3.5 Domain Theory and Propagators}
\label{sec:orgee70dd7}

Scott domains and dcpos (directed-complete partial orders) provide the
denotational semantics that connects to propagator computation:

\begin{itemize}
\item A \emph{dcpo} is a partial order where every directed set has a supremum
\item A \emph{Scott-continuous} function preserves directed suprema
\item The \emph{information ordering} on a dcpo models increasing knowledge
\end{itemize}

Propagator cells are elements of a dcpo (or complete lattice). The cell's
value increases monotonically, representing accumulating information. The
merge operation (join) is the directed supremum. Propagators are
Scott-continuous functions between dcpos.

This framing connects propagator networks to the rich theory of domain
equations, powerdomains, and continuous function spaces. In particular:
\begin{itemize}
\item \emph{Powerdomain} constructions (Plotkin, Smyth, Hoare) for nondeterminism
correspond to the three ways propagator networks can handle nondeterministic
cells (must-may analysis)
\item \emph{Bilimit-compact categories} (Amadio-Curien) ensure that recursive domain
equations have solutions --- relevant for recursive session types interpreted
as elements of recursive domains
\end{itemize}
\subsection{3.6 Coinduction and Greatest Fixpoints}
\label{sec:orgb82fb57}

Greatest fixpoints are mathematically dual to least fixpoints but operationally
different:
\begin{itemize}
\item Least fixpoint: finite proof/derivation, inductive definition
\item Greatest fixpoint: infinite behavior, coinductive definition
\end{itemize}

The canonical example is \emph{bisimulation}: two processes are bisimilar iff they
are related by the \emph{greatest fixpoint} of the bisimulation functional. This is
the coinductive definition: assume the processes are related, and verify that
every step preserves the relation.

\textbf{Coinductive types in type theory}: Coinductive types (streams, conatural
numbers, infinite trees) represent infinite data. Abel, Pientka, Thibodeau,
and Setzer (POPL 2013) introduced \emph{copatterns} for programming with
coinductive types: instead of pattern matching on constructors (inductive),
copattern matching defines behavior in response to destructors (coinductive).

\textbf{Productivity checking}: For coinductive definitions to be well-defined, they
must be \emph{productive} --- each observation step must terminate. This is the
dual of termination checking for inductive definitions.

\textbf{For Prologos}: Recursive session types with \texttt{rec} are naturally coinductive
when the protocol loops forever (a server that perpetually accepts requests).
The distinction between \(\mu\) X (inductive, eventually terminating) and \(\nu\) X
(coinductive, potentially infinite) session types maps directly to the
lfp/gfp distinction in propagator cells. A coinductive session type should be
checked for \emph{productivity} (each protocol step makes progress) rather than
\emph{termination} (the protocol eventually ends).
\subsection{3.7 Fixpoint Games on Continuous Lattices}
\label{sec:org18005ff}

Baldan, König, Padoan, and Tsampas (POPL 2019) introduced \emph{fixpoint games on
continuous lattices} --- a game-theoretic characterization of lattice fixpoints
that unifies parity games and abstract interpretation:

Given a system of fixpoint equations x\textsubscript{i} = f\textsubscript{i}(x\textsubscript{1}, \ldots{}, x\textsubscript{n}) over a
continuous lattice, they define a two-player game where:
\begin{itemize}
\item Player \(\exists\) (Existential) chooses witnesses for joins
\item Player \(\forall\) (Universal) chooses witnesses for meets
\item The winner is determined by a parity condition on the fixpoint types (\(\mu\)/\(\nu\))
\end{itemize}

\textbf{Key results}:
\begin{itemize}
\item The game characterizes the solution of the fixpoint equation system
\item For finite lattices, the game reduces to classical parity games
\item For infinite lattices, the game provides a sound approximation framework
\end{itemize}

This directly addresses the question of how propagator networks compute
nested fixpoints: the scheduler can be viewed as \emph{playing the fixpoint game},
with ascending cell updates as Existential moves and descending cell updates
as Universal moves.
\subsection{3.8 Widening and Narrowing}
\label{sec:org82c40f8}

Cousot and Cousot (1977, 1992) introduced widening (\(\nabla\)) and narrowing
(\(\Delta\)) operators to ensure termination of fixpoint iteration over
infinite-height lattices:

\begin{itemize}
\item \emph{Widening} \(\nabla\) : L \texttimes{} L \(\to\) L accelerates ascending chains by
jumping ahead, guaranteeing termination at the cost of precision (the
result is an over-approximation of the lfp)
\item \emph{Narrowing} \(\Delta\) : L \texttimes{} L \(\to\) L improves precision by descending from
the widened result, without losing the termination guarantee
\end{itemize}

Standard variants include:
\begin{itemize}
\item \emph{Delayed widening}: apply standard iteration for k steps before widening
\item \emph{Threshold widening}: widen to predefined thresholds (e.g., 0, 1, \(\infty\))
\item \emph{Guided widening}: use the property being checked to guide the widening
\end{itemize}

\textbf{For Prologos}: Widening is already implemented via the \texttt{Widenable} trait and
\texttt{net-add-widening-propagator}. The connection to abstract interpretation
suggests widening strategies informed by the temporal property being verified:
if checking AF(x < 10), widen the interval lattice with threshold 10.
\section{4. Runtime Verification and Monitoring}
\label{sec:orgf24867d}

\subsection{4.1 Three-Valued Semantics (LTL3)}
\label{sec:org4565a4e}

Runtime verification faces a fundamental challenge: LTL is defined over
infinite traces, but monitors observe only finite (growing) prefixes. Bauer,
Leucker, and Schallhart (TOSEM 2011) introduced LTL3 with three values:

\begin{itemize}
\item \textbf{T} (true): every possible infinite continuation satisfies the property
\item \textbf{F} (false): every possible infinite continuation violates the property
\item \textbf{?} (inconclusive): some continuations satisfy, some violate
\end{itemize}

Formally: [u \models \(\phi\)]\textsubscript{3} = T iff \(\forall\) w \(\in\) \(\Sigma\)\^{}\(\omega\) : u \(\cdot\) w \models \(\phi\).

The monitor construction produces a \emph{minimal deterministic finite-state machine}
that reads finite traces and outputs B3 verdicts. Extended to RV-LTL with four
values: \{T, F, T\textsubscript{p} (presumably true), F\textsubscript{p} (presumably false)\} based on
behavioral pattern extrapolation.

\textbf{Monitorable properties}: A property \(\phi\) is monitorable if for every finite
trace u, there exists a finite extension v such that u \(\cdot\) v yields a
definitive verdict. This class is \emph{strictly larger} than safety \(\cup\) co-safety.
Safety properties are finitely refutable; co-safety are finitely verifiable;
liveness properties are generally not finitely monitorable.

\textbf{Monitorable fragments of the \(\mu\)-calculus}: The alternation-free fragment (no
mutual recursion between \(\mu\) and \(\nu\)) yields monitorable properties. Research
on Hennessy-Milner logic with recursion (a \(\mu\)-calculus variant) has
established an expressiveness hierarchy of monitorable fragments.
\subsection{4.2 Monitoring Tools and Architectural Patterns}
\label{sec:orgcc07987}

The runtime verification community has developed a rich ecosystem of tools:

\textbf{Stream-based monitors}:
\begin{itemize}
\item \emph{Lola} (D'Angelo et al., TIME 2005): Declarative stream equations defining
output streams as functions of inputs. Bounded memory for syntactically
well-behaved specifications.
\item \emph{RTLola} (Finkbeiner et al., CAV 2020): Real-time extension with
asynchronous timestamped streams and temporal window aggregation. Compiles to
FPGA for hardware monitoring. Deployed on DLR ARTIS autonomous rotorcraft.
\item \emph{TeSSLa}: Non-synchronized real-time streams with recursive aggregation.
\end{itemize}

\textbf{Embedded/hardware monitors}:
\begin{itemize}
\item \emph{Copilot} (Pike et al., RV 2010): Haskell EDSL generating hard real-time
C99 monitors running in constant memory and constant time. Developed under
NASA contract for avionics.
\item \emph{R2U2} (Reinbacher et al., FMSD 2017): FPGA-based monitoring with LTL/MLTL
observers and integrated Bayesian health diagnosis. Zero software overhead.
\end{itemize}

\textbf{Parametric monitors}:
\begin{itemize}
\item \emph{JavaMOP} (Jin et al., ICSE 2012): Monitoring-Oriented Programming with
AspectJ code generation. Supports parametric properties relating specific
object instances.
\item \emph{MarQ} (Reger et al., TACAS 2015): Quantified Event Automata for
parametric monitoring with efficient trace slicing.
\end{itemize}

\begin{center}
\begin{tabular}{lll}
Pattern & Description & Examples\\
\hline
Inline & Monitor woven into system code & JavaMOP (AspectJ)\\
Sidecar & Monitor observes traces externally & Copilot, R2U2\\
Stream pipeline & Declarative stream equations & Lola, RTLola, TeSSLa\\
Hardware parallel & FPGA-based constant-time monitoring & R2U2, RTLola-FPGA\\
\end{tabular}
\end{center}

\textbf{For Prologos}: The stream-based monitoring pattern maps directly to propagator
networks. Lola's stream equations \emph{are} propagator wiring: each output stream
is a cell whose value is determined by propagators reading input streams. The
"compile temporal spec to monitor" pipeline is the same as "compile temporal
spec to propagator network." The difference is that Prologos monitors would
run at \emph{type-checking time} (static analysis) or at \emph{runtime} (dynamic
monitoring), using the same propagator infrastructure for both.
\subsection{4.3 Distributed Runtime Verification}
\label{sec:org73e99d0}

Monitoring distributed systems introduces challenges from partial ordering
and concurrent observation:

\textbf{Lamport timestamps and vector clocks}: Establish partial ordering of events
across processes. Vector clocks (Fidge 1988, Mattern 1989) precisely
characterize the happens-before relation.

\textbf{Chase-Garg lattice of consistent cuts} (1995): The set of consistent global
states forms a finite distributive lattice under inclusion. Predicate
detection on this lattice ranges from O(n\textsuperscript{2}) for linear predicates
(meet-closed) to NP-complete for general predicates.

\textbf{Decentralized monitoring} (Bauer-Falcone, FMSD 2016): Decomposing LTL
specifications into sub-formulae distributed among local monitors with
partial observation.

\textbf{For Prologos}: Distributed session types (multi-party protocols) require
distributed verification. The lattice of consistent cuts is itself a complete
lattice, and predicate detection on it is a fixpoint computation --- naturally
amenable to propagator-based evaluation. Prologos's multi-party session types
could embed distributed monitors that check protocol compliance across
endpoints.
\subsection{4.4 Predictive Runtime Verification}
\label{sec:orga7a2d92}

Predictive RV detects \emph{potential} violations that haven't occurred yet by
analyzing causal models:
\begin{itemize}
\item \emph{Happens-before analysis}: Explore alternative interleavings consistent with
the observed partial order
\item \emph{Maximal causal models}: Consider all causally compatible executions
\item \emph{SMT-based approaches}: Encode the space of possible interleavings as SMT
constraints and check for violations
\end{itemize}

This connects to the ATMS: each consistent worldview is an alternative
interleaving, and checking a property across all worldviews is predictive
verification.
\subsection{4.5 Shield Synthesis}
\label{sec:org2fcbf00}

Bloem, Könighofer, and Ehlers introduced \emph{shields} --- reactive systems that
enforce safety by observing the system's output and correcting unsafe actions:

\begin{itemize}
\item \emph{Pre-shields}: Intercept controller inputs, restricting available actions
to those guaranteed safe
\item \emph{Post-shields}: Intercept controller outputs, correcting unsafe actions to
safe alternatives
\item Formalized as safety games between the shield and the environment
\end{itemize}

\textbf{Safe reinforcement learning through shielding}: The shield constrains the
RL agent's action space, preventing unsafe exploration while maintaining
learning effectiveness. The shield is synthesized from a temporal specification
and a model of the environment.

\textbf{For Prologos}: Shield synthesis from session types would produce runtime
monitors that \emph{enforce} protocol compliance, not just detect violations. A
session-typed channel with a shield would automatically correct protocol
deviations --- a form of self-healing communication. This extends the static
verification story with dynamic enforcement.
\section{5. Session Types, Linear Logic, and Verification}
\label{sec:orgd91586f}

\subsection{5.1 The Curry-Howard Correspondence for Sessions}
\label{sec:org26d4f1f}

The Caires-Pfenning correspondence (CONCUR 2010) and Wadler's "Propositions as
Sessions" (ICFP 2012) establish a deep isomorphism between classical linear
logic and session types:

\begin{center}
\begin{tabular}{ll}
Linear Logic & Session Type\\
\hline
A \(\otimes\) B & Send A then continue as B\\
A \parr B & Receive A then continue as B\\
A \(\oplus\) B & Internal choice (select A or B)\\
A $\backslash$& B & External choice (offer A or B)\\
!A & Server (replicated service)\\
?A & Client (replicated usage)\\
Cut elimination & Communication (\(\beta\)-reduction)\\
\end{tabular}
\end{center}

\textbf{Key theorem}: Deadlock freedom follows from cut elimination in linear logic.
Well-typed processes do not deadlock because their typing derivations correspond
to valid linear logic proofs, and cut elimination preserves validity.
\subsection{5.2 Recursive Session Types: \(\mu\) vs \(\nu\)}
\label{sec:org3b4aec7}

Recursive session types come in two flavors, mirroring the lfp/gfp distinction:

\begin{itemize}
\item \textbf{Inductive} (\(\mu\)): The protocol eventually terminates. The recursive type
\(\mu\) X. !Msg. ?Ack. X describes a bounded interaction that must reach a
base case. Checked by termination/productivity analysis.

\item \textbf{Coinductive} (\(\nu\)): The protocol may run forever. The recursive type
\(\nu\) X. !Req. ?Resp. X describes a persistent server that perpetually accepts
requests. Checked by productivity: each protocol step is responsive.
\end{itemize}

The distinction is operationally critical: inductive session types correspond
to clients (finite interaction); coinductive session types correspond to
servers (persistent service). Prologos's \texttt{rec} keyword currently treats all
recursion as potentially infinite; distinguishing \(\mu\)-rec from \(\nu\)-rec would
enable more precise verification.
\subsection{5.3 Multiparty Session Types}
\label{sec:org079acd3}

Honda, Yoshida, and Carbone (POPL 2008) extended binary session types to
multi-party protocols via \emph{global types} --- protocol specifications describing
the interaction among all participants:

\begin{enumerate}
\item \emph{Global type}: Describes the protocol from a bird's-eye view
\item \emph{Projection}: Each participant's local type is projected from the global type
\item \emph{Deadlock freedom}: If all participants follow their projected local types,
the system is deadlock-free
\end{enumerate}

Recent extensions include:
\begin{itemize}
\item \emph{Timed multiparty session types} (ECOOP 2024): Adding timing constraints
\item \emph{Crash-stop multiparty session types} (ECOOP 2023): Handling participant
failures
\item \emph{Behavioral contracts} (Castagna-Gesbert-Padovani 2009): Subtyping for
session types ensuring safe substitution
\item \emph{Gradual session types} (Igarashi et al. 2017): Blending static and
dynamic session type checking
\end{itemize}
\subsection{5.4 Session Types and Model Checking}
\label{sec:org0774c5a}

Lange, Tuosto, and Yoshida developed Communicating Session Automata (CSA) ---
a direct bridge between session types and model checking:

\begin{itemize}
\item Global types are translated to communicating finite-state machines (CFSMs)
\item Properties (deadlock freedom, protocol conformance) are checked via model
checking on the product CFSM
\item The model checking can verify temporal properties beyond what type checking
alone can ensure
\end{itemize}

This is precisely the bridge Prologos aims to build, but using propagator
networks instead of CFSMs as the verification substrate.
\subsection{5.5 Actris: Session Types Meet Separation Logic}
\label{sec:org6cc348d}

Hinrichsen, Bengtson, and Krebbers (POPL 2020) introduced Actris ---
\emph{Dependent Separation Protocols} within the Iris framework. This merges session
types with higher-order separation logic, where protocol steps can depend on
separation logic propositions about heap state.

Actris 2.0 (LMCS 2022) adds subprotocols (session-type subtyping) for
compositional reasoning. This is the most sophisticated integration of
session types with formal verification to date, and demonstrates that
session types and separation logic are not competing but complementary
verification approaches.
\section{6. Formal Verification Frontiers}
\label{sec:org3507473}

\subsection{6.1 Separation Logic and Iris}
\label{sec:orgd124c6f}

Iris (Jung et al., POPL 2015) is a higher-order concurrent separation logic
framework mechanized in Coq, built on three orthogonal pillars:

\begin{enumerate}
\item \textbf{Resource Algebras (Cameras)}: Generalized partial commutative monoids with
a validity predicate --- the algebraic foundation for resource reasoning
\item \textbf{Invariants}: Named propositions that hold at all times, with impredicative
invariant opening
\item \textbf{Ghost State}: Logical state not in the physical program, governed by
resource algebras and frame-preserving updates
\end{enumerate}

The \emph{later modality} (\triangleright P) resolves circularity in impredicative
definitions via guarded recursion.

\textbf{RustBelt} (Jung et al., POPL 2018) proves Rust's type system sound, including
unsafe code, using Iris's lifetime logic. \textbf{RefinedRust} (Sammler et al., PLDI
\begin{enumerate}
\item layers refinement types on Rust within Iris, with machine-checked proofs.
\end{enumerate}
\subsection{6.2 Refinement Types}
\label{sec:org0b467d2}

\textbf{Liquid types} (Rondon-Kawaguchi-Jhala, PLDI 2008) combine HM type inference
with predicate abstraction: types \{v : T | p\} where p is drawn from a decidable
qualifier language. Subtyping reduces to SMT implication checking.

\textbf{Refinement reflection} (Vazou et al., POPL 2018) enables complete verification
by reflecting function definitions into their output refinement types.

\textbf{F*} (Swamy et al.) combines refinement types with an \emph{indexed effect system}
organized as a lattice of effects, enabling modular verification of stateful
programs.

\textbf{For Prologos}: Refinement types could extend Prologos's dependent types with
SMT-backed lightweight verification. The qualifier lattice in liquid type
inference is a complete lattice amenable to propagator computation ---
qualifier inference \emph{is} fixpoint computation over the qualifier lattice.
\subsection{6.3 Algebraic Effects and Handlers}
\label{sec:orge5949ad}

Algebraic effects (Plotkin-Power 2001-2003) model computational effects as
operations of an algebraic theory. Effect handlers (Plotkin-Pretnar, ESOP 2009)
generalize exception handlers with access to delimited continuations.

Key languages: Eff, Frank (where handlers are the universal abstraction), Koka
(row-polymorphic effect types compiling to C). Recent work (Sekiyama-Tsukada,
POPL 2024) combines refinement types with algebraic effects.

\textbf{Connection to session types}: Communication operations (send, receive) can be
modeled as algebraic effects; session types then become effect types constraining
operation ordering. Both algebraic effects (via free monads) and session types
(via indexed monads) unify under graded monadic frameworks.
\subsection{6.4 Reactive Synthesis}
\label{sec:orge444803}

Given a temporal specification, \emph{synthesize} an implementation automatically:

\begin{itemize}
\item \textbf{Church's problem} (1957): Realized by the Büchi-Landweber theorem (1969)
for \(\omega\)-regular specifications
\item \textbf{GR(1) synthesis} (Bloem-Jobstmann-Piterman-Pnueli-Sa'ar, JCSS 2012):
Efficient O(n \(\cdot\) m \(\cdot\) N\textsuperscript{2}) synthesis for the practically important
generalized reactivity(1) fragment
\item \textbf{LTL synthesis}: 2EXPTIME-complete (Pnueli-Rosner 1989); Safraless
approaches (Kupferman-Vardi, FOCS 2005) avoid determinization
\item \textbf{Bounded synthesis} (Finkbeiner-Schewe, ATVA 2007): Restrict to
implementations with at most k states; reduces to SAT
\end{itemize}

\textbf{For Prologos}: Reactive synthesis from session type specifications is a
natural extension of the propagator-based verification story. Given temporal
properties on session types, synthesize process implementations that satisfy
them. The propagator network supports this through its \emph{bidirectional} nature
--- constraints flow in both directions, enabling constructive derivation of
implementations from specifications.
\subsection{6.5 Coalgebraic Semantics}
\label{sec:orgd494936}

Coalgebras provide a unified categorical treatment of transition systems and
temporal logics:

\begin{itemize}
\item A coalgebra for functor F is a morphism c : X \(\to\) F(X)
\item Different F capture different system types: P(X) for nondeterminism,
D(X) for probability, (2 \texttimes{} X)\textsuperscript{A} for deterministic automata
\item Coalgebraic modal logic uses predicate liftings to define modal operators
parametrically in the functor
\end{itemize}

This enables \emph{parametric verification}: fixing the functor determines the
system type, and results transfer across functor choices. Stone duality
connects coalgebraic logic to topology, yielding representation theorems
and completeness results.

\textbf{For Prologos}: Prologos's propagator network is itself a coalgebra ---
the network state transforms via propagator application. Temporal properties
of the network (convergence, stability, monotonicity) can be expressed in
coalgebraic modal logic. The terminal coalgebra (greatest fixpoint) of the
network functor represents the "fully resolved" type/session/effect
assignment.
\subsection{6.6 Capability-Based Security Verification}
\label{sec:orgb6c28d9}

Object-capability (oCap) security restricts authority propagation to explicit
capability passing. Formal verification includes:

\begin{itemize}
\item \textbf{Cerise} (Georges et al., JACM 2023): Verifying capability machine programs
using Iris, proving security even with untrusted code
\item \textbf{CHERI} (Watson et al.): Hardware-enforced capabilities with unforgeable
tagged pointers; VeriCHERI provides RTL-level verification
\item \textbf{WebAssembly}: Capability-based module system with explicit import/export
surfaces
\end{itemize}

\textbf{Connection to linear types}: Capability security and linear types share deep
structure --- both enforce resource discipline. Affine capabilities (usable at
most once) are linear resources. CHERI's hardware tag enforces a form of
linearity. Cerise explicitly uses Iris's resource algebras (substructural
reasoning) to verify capability programs.

\textbf{For Prologos}: Prologos's linear types (QTT) already provide resource control;
adding a capability-based module system would formalize authority propagation.
The propagator network can verify capability confinement properties: a
capability that should not escape a scope is a linear resource that must not
be sent outside its boundary --- checkable via the session type checker.
\section{7. The Propagator Network as Universal Verification Substrate}
\label{sec:org1bdabd5}

\subsection{7.1 The Fundamental Observation}
\label{sec:orgdfa6074}

The preceding survey reveals a common mathematical structure across all
formalisms:

\textbf{Every verification formalism discussed in this document computes fixpoints
over complete lattices.}

\begin{center}
\begin{tabular}{lll}
Formalism & Lattice & Fixpoint Type\\
\hline
CTL model checking & (2\textsuperscript{S}, \subseteq) & lfp and gfp\\
\(\mu\)-calculus & (2\textsuperscript{S}, \subseteq) & nested lfp/gfp\\
STL robustness & (\overline\{\mathbb{R}\}, \(\le\)) & inf/sup over time\\
Abstract interpretation & Abstract domain via Galois conn. & lfp with widening\\
AFT / bilattice & L\textsuperscript{2} with precision ordering & Kripke-Kleene/stable\\
PCTL probabilistic & ([0,1], \(\le\)) & lfp via lin. eqs.\\
Energy games & (\mathbb{Z}, \(\le\)) & lfp/gfp combined\\
Type inference & Type lattice & lfp via unification\\
Session type checking & Session lattice & lfp (currently)\\
LTL3 monitoring & (\{T, F, ?\}, \(\le\)\textsubscript{k}) & lfp (trace prefixes)\\
Liquid type inference & Qualifier lattice & lfp via abstraction\\
Reactive synthesis & Game lattice & nested lfp/gfp\\
\end{tabular}
\end{center}

Propagator networks are a \emph{general-purpose engine} for lattice fixpoint
computation. The network infrastructure --- cells, propagators, scheduler,
ATMS --- is agnostic to the specific lattice. Adding a new verification
capability means:
\begin{enumerate}
\item Define the lattice (implement the \texttt{Lattice} trait)
\item Define the propagators (monotone functions between cells)
\item Wire them into the existing network
\item \texttt{run-to-quiescence} computes the fixpoint
\end{enumerate}
\subsection{7.2 What Prologos Already Has (Infrastructure Inventory)}
\label{sec:orgbd3a5f2}

From \texttt{propagator.rkt} and the existing verification layers:

\begin{center}
\begin{tabular}{lll}
Component & What It Provides & Verification Role\\
\hline
Cells (ascending) & Lattice-valued mutable state & State labeling, property cells\\
Monotone propagators & Functions between cells & Transition relation, operators\\
BSP/Jacobi scheduler & Parallel fixpoint iteration & Model checking algorithm\\
ATMS & Multi-hypothesis reasoning & Path quantification (A/E)\\
Nogoods & Inconsistent assumption sets & Impossible paths / counterexamples\\
Galois connections & Cross-domain abstraction & Abstract model checking\\
Widening (\texttt{Widenable}) & Accelerated convergence & Termination for infinite lattices\\
Threshold propagators & Wait-for-value patterns & Stratified evaluation barriers\\
Cross-domain propagators & Inter-lattice bridges & Multi-domain verification\\
\end{tabular}
\end{center}
\subsection{7.3 What Would Complete the Picture}
\label{sec:orgc48f1a0}

\subsubsection{7.3.1 Greatest Fixpoint Support (Descending Cells)}
\label{sec:org269bcfb}

Currently all cells start at \bot and refine upward (lfp). Safety properties
and coinductive definitions require the dual: cells starting at \top and
refining downward (gfp).

\textbf{Implementation sketch}:
\begin{enumerate}
\item Add \texttt{cell-direction} flag: \texttt{:ascending} (default) or \texttt{:descending}
\item Descending cells use \texttt{meet} instead of \texttt{join} for \texttt{net-cell-write}
\item Contradiction for descending cells is reaching \bot (dual of ascending \top)
\item The scheduler processes both directions; nested fixpoints stabilize
inside-out
\end{enumerate}

This is architecturally straightforward. The session lattice already has both
join and meet. The extension adds the dual direction to the scheduler.
\subsubsection{7.3.2 Temporal Property Specification Language}
\label{sec:org25444a1}

A user-facing language for specifying temporal properties that compiles to
propagator wiring (developed in detail in \texttt{PROPAGATORS\_AS\_MODEL\_CHECKERS.org}
§4):

\begin{verbatim}
property pin-safe
  :session Pin-Guard
  :always [not [and :pinned :gc-collecting]]

property request-response
  :session CapTP-Steady
  :always [implies [sent :call] [eventually [received :return]]]
\end{verbatim}

Each keyword maps to a fixpoint construction:
\begin{itemize}
\item \texttt{:always} \(\to\) \(\nu\) X. \(\phi\) \(\land\) {[}a]X (greatest fixpoint)
\item \texttt{:eventually} \(\to\) \(\mu\) X. \(\phi\) \(\lor\) {[}a]X (least fixpoint)
\item \texttt{:until} \(\to\) \(\mu\) X. \(\psi\) \(\lor\) (\(\phi\) \(\land\) {[}a]X) (least fixpoint)
\end{itemize}
\subsubsection{7.3.3 Approximation Fixpoint Theory Integration}
\label{sec:org244ceaf}

Using AFT's bilattice construction L\textsuperscript{2} for handling non-monotone operations:
\begin{itemize}
\item Lower bound cell (what is certainly true) + upper bound cell (what is
possibly true)
\item Stable operator applied to the pair
\item Generalizes the stratification controller to arbitrary non-monotone
operators
\end{itemize}
\subsubsection{7.3.4 Runtime Monitoring Compilation}
\label{sec:org02781cd}

Temporal properties that cannot be verified statically (due to undecidability
or state-space explosion) compile to \emph{runtime monitors} --- propagator networks
that run alongside the program:
\begin{itemize}
\item Stream-based monitoring (Lola-style) maps directly to propagator wiring
\item The same temporal specification language works for both static and dynamic
verification
\item Shields (enforcement monitors) correct unsafe behavior rather than just
detecting it
\end{itemize}
\subsubsection{7.3.5 Quantitative Verification Support}
\label{sec:orgbc47f4a}

Replace the Boolean lattice with richer domains:
\begin{itemize}
\item (\overline\{\mathbb{R}\}, \(\le\)) for STL robustness monitoring
\item ([0,1], \(\le\)) for probabilistic verification
\item (\mathbb{Z}, \(\le\)) for energy/resource constraints
\end{itemize}

The propagator scheduler is lattice-agnostic; only the cell merge function
changes.
\subsection{7.4 The Verification Layer Cake (Extended)}
\label{sec:org30815ae}

\begin{verbatim}
Layer 7: Hyperproperties          (HyperLTL: information flow, non-interference)
     |                              → ATMS worldview quantification
Layer 6: Temporal Properties      (CTL/LTL/μ-calculus: safety, liveness, fairness)
     |                              → Fixpoint propagators (ascending + descending)
Layer 5: Runtime Monitors         (Stream-based monitoring, shields)
     |                              → Compiled propagator sub-networks
Layer 4: Effect Ordering          (Causal model: deadlock detection)
     |                              → Galois connection to session domain
Layer 3: Session Types            (Protocol state machines: duality, completeness)
     |                              → Session lattice, decomposition propagators
Layer 2: Linear Types (QTT)       (Resource tracking: use-once, capabilities)
     |                              → Multiplicity lattice, QTT propagators
Layer 1: Dependent Types          (Value-dependent properties: indexed types)
     |                              → Type lattice, unification propagators
Layer 0: Propagator Network       (Monotone fixpoint computation: universal substrate)
         + ATMS (path quantification)
         + Galois Connections (abstraction)
         + Widening (convergence acceleration)
         + AFT/Bilattice (non-monotone recovery)
\end{verbatim}

Every layer builds on Layer 0. Every new layer adds cells and propagators to
the \emph{same network}. The cross-layer interactions happen through the shared
substrate. A temporal property that references type information, session state,
effect ordering, and resource tracking is verified in a single
\texttt{run-to-quiescence} pass because all layers share the network.
\section{8. Cross-Cutting Themes}
\label{sec:org31d1549}

\subsection{8.1 The Lattice Universality Principle}
\label{sec:org39ba87f}

The recurring theme across all formalisms is that \emph{verification is lattice
computation}. Different formalisms choose different lattices and different
fixpoint operators, but the computational pattern is invariant:

\begin{enumerate}
\item Define a complete lattice
\item Define monotone functions over it
\item Compute fixpoints (least, greatest, or nested)
\item Interpret the fixpoint as a verification result
\end{enumerate}

Propagator networks are the most general realization of this pattern: they are
\emph{parameterized} by the lattice and the functions, and the scheduler computes
fixpoints for \emph{any} instantiation.
\subsection{8.2 The Duality of Safety and Liveness}
\label{sec:org2c4d014}

Safety (\emph{nothing bad happens}) and liveness (\emph{something good eventually
happens}) are the two fundamental verification concerns. They correspond
exactly to the two fixpoint types:

\begin{center}
\begin{tabular}{llll}
Concern & Fixpoint & Cell Direction & Iteration\\
\hline
Safety & gfp (\(\nu\)) & Descending & Start at \top, refine down\\
Liveness & lfp (\(\mu\)) & Ascending & Start at \bot, refine up\\
\end{tabular}
\end{center}

The Alpern-Schneider decomposition theorem (1985) states that every temporal
property is the intersection of a safety and a liveness property. In
propagator terms: every temporal property can be decomposed into an ascending
cell (liveness component) and a descending cell (safety component), with
propagators connecting them.
\subsection{8.3 Static vs. Dynamic Verification}
\label{sec:org93541a4}

The same propagator infrastructure supports both static (compile-time) and
dynamic (runtime) verification:

\begin{center}
\begin{tabular}{lll}
Aspect & Static Verification & Dynamic Verification\\
\hline
Input & Type/session annotations & Runtime events/signals\\
Lattice & Type/session lattice & Trace prefix lattice / [0,1]\\
Completeness & May be incomplete & Always reflects actual behavior\\
Overhead & Compile-time only & Runtime performance cost\\
Expressiveness & Decidable properties only & Any monitorable property\\
Implementation & Propagators in type checker & Propagators in runtime system\\
\end{tabular}
\end{center}

The key insight is that \emph{the propagator wiring is the same}. A temporal
property specification compiles to the same propagator structure whether it
runs during type checking or during execution. The difference is only in
the source of cell updates: type annotations vs. runtime events.
\subsection{8.4 The Abstraction-Refinement Paradigm}
\label{sec:orgdd0b2ff}

Abstract interpretation, CEGAR, and Galois connections are all manifestations
of the same pattern: compute on a simpler (abstract) domain, and refine when
the abstraction is too coarse.

In propagator terms:
\begin{enumerate}
\item Create abstract cells connected to concrete cells via Galois connection
propagators
\item Compute the fixpoint on the abstract domain (faster, may be imprecise)
\item If the abstract result is inconclusive, refine the abstraction (add finer
abstract cells) and re-run
\item The propagator network handles this incrementally --- only affected cells
are re-evaluated
\end{enumerate}
\subsection{8.5 Resource Discipline as Universal Theme}
\label{sec:org566fd82}

Across the formalisms, \emph{resource discipline} --- controlling who can use what,
when, and how many times --- is the universal verification theme:

\begin{itemize}
\item \textbf{Linear types (QTT)}: Resources used exactly once (or controlled multiplicity)
\item \textbf{Session types}: Communication channels used according to protocol
\item \textbf{Capabilities}: Authorities that cannot be forged
\item \textbf{Energy constraints}: Resources consumed/produced along execution
\item \textbf{Separation logic}: Heap resources owned exclusively or shared under invariant
\item \textbf{Algebraic effects}: Computational effects requiring explicit handling
\end{itemize}

All enforce the same principle at different granularities. Prologos's
combination of QTT + session types + capability security already covers the
first three; extending to energy constraints and separation-logic reasoning
would complete the resource verification picture.
\section{9. Research Program}
\label{sec:org0b52054}

\subsection{9.1 Immediate (Infrastructure Extensions)}
\label{sec:orgdd41f6b}

\begin{enumerate}
\item \textbf{Greatest fixpoint cells}: Add \texttt{:descending} cell direction to
\texttt{propagator.rkt}. Validate with bisimulation checking as the first
application.

\item \textbf{Mixed ascending/descending networks}: Verify that the BSP scheduler
correctly handles mixed-direction cells with nested fixpoint computation.
Test with Level 2 alternation (\(\nu\) X. \(\mu\) Y. \(\phi\)).

\item \textbf{STL robustness cells}: Implement
\overline\{\mathbb{R}\}-valued cells with inf/sup merge. Validate with
simple continuous signal monitoring.
\end{enumerate}
\subsection{9.2 Medium-Term (New Verification Capabilities)}
\label{sec:org936b696}

\begin{enumerate}
\item \textbf{Temporal property language}: Design and implement the \texttt{property} keyword
with temporal operators (:always, :eventually, :until, :within). Compile to
propagator wiring.

\item \textbf{AFT integration}: Implement bilattice L\textsuperscript{2} cell pairs for non-monotone
operator recovery. Generalize the stratification controller.

\item \textbf{Runtime monitor compilation}: Compile temporal specifications to Lola-style
stream monitors that can run alongside Prologos programs.

\item \textbf{Multiparty session type verification}: Extend session type propagators to
handle global types with projection and multi-participant deadlock checking.
\end{enumerate}
\subsection{9.3 Long-Term (Research Frontiers)}
\label{sec:org41d061b}

\begin{enumerate}
\item \textbf{Reactive synthesis from session types}: Given temporal properties on a
session type, synthesize conforming process implementations. Leverage the
bidirectional nature of propagator networks.

\item \textbf{Probabilistic session types}: Session types with probabilistic branching
and PCTL-style probability bounds on protocol properties.

\item \textbf{Hyperproperty verification}: Leverage ATMS worldview quantification
to verify HyperLTL properties of session-typed protocols (information
flow, non-interference).

\item \textbf{Coinductive session types}: Distinguish \(\mu\)-rec (inductive/terminating)
from \(\nu\)-rec (coinductive/persistent) and check productivity rather than
termination for persistent servers.

\item \textbf{Capability confinement verification}: Verify that capabilities do not
escape their intended scope, using propagator-based reachability analysis
on the session type's communication graph.

\item \textbf{Shield synthesis for session types}: Synthesize runtime enforcement
shields from session type specifications, producing self-healing
communication channels.
\end{enumerate}
\section{10. References}
\label{sec:org9466ad8}

\subsection{Modal \(\mu\)-Calculus and Fixpoint Theory}
\label{sec:org3a1d80d}
\begin{itemize}
\item Kozen, D. "Results on the propositional \(\mu\)-calculus." \emph{TCS} 27(3), 1983.
\item Walukiewicz, I. "Completeness of Kozen's axiomatisation." \emph{Information and Computation} 157(1-2), 2000.
\item Janin, D. and Walukiewicz, I. "On the expressive completeness of the propositional \(\mu\)-calculus with respect to MSO." \emph{CONCUR} 1996.
\item Bradfield, J. "The modal \(\mu\)-calculus alternation hierarchy is strict." \emph{CONCUR} 1996.
\item Calude, C., Jain, S., Khoussainov, B., Li, W., and Stephan, F. "Deciding parity games in quasipolynomial time." \emph{STOC} 2017.
\item Jurdzinski, M. "Small progress measures for solving parity games." \emph{STACS} 2000.
\item Jurdzinski, M. and Lazic, R. "Succinct progress measures for solving parity games." \emph{LICS} 2017.
\item Baldan, P., König, B., Padoan, T., and Tsampas, C. "Fixpoint games on continuous lattices." \emph{POPL} 2019.
\item Emerson, E.A. and Jutla, C. "Tree automata, mu-calculus and determinacy." \emph{FOCS} 1991.
\end{itemize}
\subsection{Temporal Logics}
\label{sec:org037c2d6}
\begin{itemize}
\item Pnueli, A. "The temporal logic of programs." \emph{FOCS} 1977.
\item Clarke, E., Emerson, E., and Sistla, A. "Automatic verification of finite-state concurrent systems using temporal logic specifications." \emph{ACM TOPLAS} 8(2), 1986.
\item Emerson, E. and Halpern, J. "'Sometimes' and 'not never' revisited: on branching versus linear time temporal logic." \emph{JACM} 33(1), 1986.
\item Maler, O. and Nickovic, D. "Monitoring temporal properties of continuous signals." \emph{FORMATS} 2004.
\item Fainekos, G. and Pappas, G. "Robustness of temporal logic specifications for continuous-time signals." \emph{TCS} 410(42), 2009.
\item Donzé, A. and Maler, O. "Robust satisfaction of temporal logic over real-valued signals." \emph{FORMATS} 2010.
\item Koymans, R. "Specifying real-time properties with metric temporal logic." \emph{RTS} 2(4), 1990.
\item Alur, R., Feder, T., and Henzinger, T. "The benefits of relaxing punctuality." \emph{JACM} 43(1), 1996.
\item Ouaknine, J. and Worrell, J. "On the decidability of metric temporal logic." \emph{LICS} 2005.
\end{itemize}
\subsection{Hyperproperties and Information Flow}
\label{sec:orgb643c73}
\begin{itemize}
\item Clarkson, M. and Schneider, F. "Hyperproperties." \emph{JCS} 18(6), 2010.
\item Clarkson, M., Finkbeiner, B., Koleini, M., Rabe, M., and Sanchez, C. "Temporal logics for hyperproperties." \emph{POST} 2014.
\end{itemize}
\subsection{Probabilistic and Quantitative Verification}
\label{sec:org1ea8744}
\begin{itemize}
\item Hansson, H. and Jonsson, B. "A logic for reasoning about time and reliability." \emph{Formal Aspects of Computing} 6(5), 1994.
\item Baier, C., Haverkort, B., Hermanns, H., and Katoen, J.-P. "Model-checking algorithms for continuous-time Markov chains." \emph{IEEE TSE} 29(6), 2003.
\item Kwiatkowska, M., Norman, G., and Parker, D. "PRISM 4.0: Verification of probabilistic real-time systems." \emph{CAV} 2011.
\item Almagor, S., Boker, U., and Kupferman, O. "Discounting in LTL." \emph{TACAS} 2014.
\item Chatterjee, K. and Doyen, L. "Energy parity games." \emph{ICALP} 2010.
\item Kaminski, B., Katoen, J.-P., Matheja, C., and Olmedo, F. "Weakest precondition reasoning for expected runtimes of randomized algorithms." \emph{JACM} 2018.
\end{itemize}
\subsection{Runtime Verification}
\label{sec:orgcb36ad3}
\begin{itemize}
\item Bauer, A., Leucker, M., and Schallhart, C. "Runtime verification for LTL and TLTL." \emph{ACM TOSEM} 20(4), 2011.
\item Leucker, M. and Schallhart, C. "A brief account of runtime verification." \emph{JLAP} 78(5), 2009.
\item D'Angelo, B. et al. "LOLA: Runtime monitoring of synchronous systems." \emph{TIME} 2005.
\item Finkbeiner, B. et al. "RTLola cleared for take-off: Monitoring autonomous aircraft." \emph{CAV} 2020.
\item Pike, L. et al. "Copilot: A hard real-time runtime monitor." \emph{RV} 2010.
\item Reinbacher, T. et al. "R2U2: Monitoring and diagnosis of security threats for unmanned aerial systems." \emph{FMSD} 2017.
\item Jin, D., Meredith, P., Lee, C., and Rosu, G. "JavaMOP: Efficient parametric runtime monitoring framework." \emph{ICSE} 2012.
\end{itemize}
\subsection{Approximation Fixpoint Theory and Bilattices}
\label{sec:orgd380b49}
\begin{itemize}
\item Denecker, M., Marek, V., and Truszczyński, M. "Approximations, stable operators, well-founded fixpoints and applications in nonmonotonic reasoning." \emph{Logic-Based AI} 2000.
\item Denecker, M., Marek, V., and Truszczyński, M. "Ultimate approximation and its application in nonmonotonic knowledge representation systems." \emph{Information and Computation} 192(1), 2004.
\item Belnap, N. "A useful four-valued logic." \emph{Modern Uses of Multiple-Valued Logic} 1977.
\item Ginsberg, M. "Multivalued logics: A uniform approach to reasoning in AI." \emph{Computational Intelligence} 4, 1988.
\item Fitting, M. "Bilattices and the semantics of logic programming." \emph{JLP} 11(2), 1991.
\end{itemize}
\subsection{Abstract Interpretation}
\label{sec:orgcff7f58}
\begin{itemize}
\item Cousot, P. and Cousot, R. "Abstract interpretation: A unified lattice model for static analysis." \emph{POPL} 1977.
\item Cousot, P. and Cousot, R. "Comparing the Galois connection and widening/narrowing approaches to abstract interpretation." \emph{PLILP} 1992.
\item Clarke, E., Grumberg, O., and Long, D. "Model checking and abstraction." \emph{ACM TOPLAS} 16(5), 1994.
\end{itemize}
\subsection{Session Types and Linear Logic}
\label{sec:org768b3d7}
\begin{itemize}
\item Caires, L. and Pfenning, F. "Session types as intuitionistic linear propositions." \emph{CONCUR} 2010.
\item Wadler, P. "Propositions as sessions." \emph{JFP} 24(2-3), 2014.
\item Honda, K., Yoshida, N., and Carbone, M. "Multiparty asynchronous session types." \emph{POPL} 2008.
\item Hinrichsen, J., Bengtson, J., and Krebbers, R. "Actris: Session-type based reasoning in separation logic." \emph{POPL} 2020.
\end{itemize}
\subsection{Separation Logic}
\label{sec:org978203d}
\begin{itemize}
\item Jung, R. et al. "Iris: Monoids and invariants as an orthogonal basis for concurrent reasoning." \emph{POPL} 2015.
\item Jung, R. et al. "Iris from the ground up." \emph{JFP} 28, 2018.
\item Jung, R. et al. "RustBelt: Securing the foundations of the Rust programming language." \emph{POPL} 2018.
\item Sammler, M. et al. "RefinedRust: A type system for high-assurance verification of Rust programs." \emph{PLDI} 2024.
\end{itemize}
\subsection{Refinement Types and Verification}
\label{sec:orgf037ec9}
\begin{itemize}
\item Rondon, P., Kawaguchi, M., and Jhala, R. "Liquid types." \emph{PLDI} 2008.
\item Vazou, N. et al. "Refinement reflection: Complete verification with SMT." \emph{POPL} 2018.
\item Swamy, N. et al. "Dependent types and multi-monadic effects in F*." \emph{POPL} 2016.
\end{itemize}
\subsection{Algebraic Effects}
\label{sec:org708df32}
\begin{itemize}
\item Plotkin, G. and Pretnar, M. "Handlers of algebraic effects." \emph{ESOP} 2009.
\item Lindley, S., McBride, C., and McLaughlin, C. "Do be do be do." \emph{POPL} 2017.
\end{itemize}
\subsection{Reactive Synthesis}
\label{sec:orgeab16d5}
\begin{itemize}
\item Bloem, R., Jobstmann, B., Piterman, N., Pnueli, A., and Sa'ar, Y. "Synthesis of reactive(1) designs." \emph{JCSS} 78(3), 2012.
\item Pnueli, A. and Rosner, R. "On the synthesis of a reactive module." \emph{POPL} 1989.
\item Finkbeiner, B. and Schewe, S. "Bounded synthesis." \emph{ATVA} 2007.
\end{itemize}
\subsection{Capability Security}
\label{sec:org56859bc}
\begin{itemize}
\item Swasey, D., Garg, D., and Dreyer, D. "Robust and compositional verification of object capability patterns." \emph{OOPSLA} 2017.
\item Georges, A. et al. "Cerise: Program verification on a capability machine in the presence of untrusted code." \emph{JACM} 2023.
\end{itemize}
\subsection{Coalgebraic Methods}
\label{sec:org9afd08f}
\begin{itemize}
\item Kupke, C. and Pattinson, D. "Coalgebraic semantics of modal logics: An overview." \emph{TCS} 2011.
\item Hasuo, I. "Tracing anonymity with coalgebras." PhD Thesis, Radboud University, 2008.
\end{itemize}
\subsection{Domain Theory and Coinduction}
\label{sec:org9c1a458}
\begin{itemize}
\item Abel, A., Pientka, B., Thibodeau, D., and Setzer, A. "Copatterns: Programming infinite structures by observations." \emph{POPL} 2013.
\item Amadio, R. and Curien, P.-L. \emph{Domains and Lambda-Calculi}. Cambridge, 1998.
\end{itemize}
\subsection{Propagator Networks}
\label{sec:orgf1e283b}
\begin{itemize}
\item Radul, A. "Propagation Networks: A Flexible and Expressive Substrate for Computation." PhD Thesis, MIT, 2009.
\item Radul, A. and Sussman, G. "The Art of the Propagator." MIT CSAIL, 2009.
\item De Kleer, J. "An assumption-based TMS." \emph{Artificial Intelligence} 28(2), 1986.
\end{itemize}
\subsection{Prologos Project Internal Documents}
\label{sec:org812b1eb}
\begin{itemize}
\item "Propagator Networks as Model Checkers." \texttt{docs/research/2026-03-07\_PROPAGATORS\_AS\_MODEL\_CHECKERS.org}. 2026.
\item "The Categorical Structure of Layered Recovery." \texttt{docs/research/2026-03-13\_LAYERED\_RECOVERY\_CATEGORICAL\_ANALYSIS.org}. 2026.
\item "Propagator Networks: A Comprehensive Survey." \texttt{docs/research/PROPAGATOR\_NETWORKS.org}. 2026.
\end{itemize}
\end{document}
